\begin{textsample}{POS Dim 1 – 2024 – Score 19.00 – 2024/t003234.txt}  \label{ex:f1_pos_049}
I want to tell you what I see coming.
I ' ve been lucky enough to be working on AI for almost 15 years now.
Back when I started, to describe it as fringe would be an understatement.
Researchers would say,"No, no, we ' re only working on machine \textbf{learning}."Because working on AI was seen as way too out there.
In 2010, just \textbf{the} very mention of \textbf{the} phrase"AGI,"\textbf{artificial} general intelligence, would get you some seriously strange looks and even a cold shoulder."You ' re actually building AGI?"people would say."Isn ' t that something out of science fiction?"People thought it was 50 years away or 100 years away, if it was even possible at all.
Talk of AI was, I guess, kind of embarrassing.
People generally thought we were weird.
And I guess in some ways we kind of were.
It wasn ' t long, though, before AI started beating humans at a whole range of tasks that people previously thought were way out of reach.
Understanding images, translating languages, transcribing speech, playing Go and chess and even diagnosing diseases.
People started waking up to \textbf{the} fact that AI was going to have an enormous impact, and they were rightly asking technologists like me some pretty tough questions.
Is it true that AI is going to solve \textbf{the} climate crisis?
Will it make personalized education available to everyone?
Does it mean we ' ll all get universal basic income and we won ' t have to work anymore?
Should I be afraid?
What does it mean for weapons and war?
And of course, will China win?
Are we in a race?
Are we headed for a mass misinformation apocalypse?
All good questions.
But it was actually a simpler and much more kind of fundamental question that left me puzzled.
One that actually gets to \textbf{the} very heart of my work every day.
One morning over breakfast, my six-year-old nephew Caspian was playing with Pi, \textbf{the} AI I created at my last company, Inflection.
With a mouthful of scrambled eggs, he looked at me plain in \textbf{the} face and said,"But Mustafa, what is an AI anyway?"He ' s such a sincere and curious and optimistic little guy.
He ' d been talking to Pi about how cool it would be if one day in \textbf{the} future, he could visit dinosaurs at \textbf{the} zoo.
And how he could make infinite amounts of chocolate at home.
And why Pi couldn ' t yet play I Spy."Well,"I said,"it ' s a clever piece of software that ' s read most of \textbf{the} text on \textbf{the} open internet, and it can talk to you about anything you want.""Right.
So like a person then?"I was stumped.
Genuinely left scratching my head.
All my boring stock answers came rushing through my mind."No, but AI is just another general-purpose technology, like printing or steam."It will be a tool that will augment us and make us smarter and more productive.
And when it gets better over time, it ' ll be like an all-knowing oracle that will help us solve grand scientific challenges."You know, all of these responses started to feel, I guess, a little bit defensive.
And actually better suited to a policy seminar than breakfast with a no-nonsense six-year-old."Why am I hesitating?"I thought to myself.
You know, let ' s be honest.
My nephew was asking me a simple question that those of us in AI just don ' t confront often enough.
What is it that we are actually creating?
What does it mean to make something totally new, fundamentally different to any invention that we have known before?
It is clear that we are at an inflection point in \textbf{the} history of humanity.
On our current trajectory, we ' re headed towards \textbf{the} emergence of something that we are all struggling to describe, and yet we cannot control what we don ' t understand.
And so \textbf{the} metaphors, \textbf{the} mental models, \textbf{the} names, these all matter if we ' re to get \textbf{the} most out of AI whilst limiting its potential downsides.
As someone who embraces \textbf{the} possibilities of this technology, but who ' s also always cared deeply about its ethics, we should, I think, be able to easily describe what it is we are building.
And that includes \textbf{the} six-year-olds.
So it ' s in that spirit that I offer up today \textbf{the} following metaphor for helping us to try to grapple with what this moment really is.
I think AI should best be understood as something like a new digital species.
Now, don ' t take this too literally, but I predict that we ' ll come to see them as digital companions, new partners in \textbf{the} journeys of all our lives.
Whether you think we ' re on a 10-, 20- or 30-year path here, this is, in my view, \textbf{the} most accurate and most fundamentally honest way of describing what ' s actually coming.
And above all, it enables everybody to prepare for and shape what comes next.
Now I totally get, this is a strong claim, and I ' m going to explain to everyone as best I can why I ' m making it.
But first, let me just try to set \textbf{the} context.
From \textbf{the} very first microscopic organisms, life on Earth stretches back billions of years.
Over that time, life \textbf{evolved} and diversified.
Then a few million years ago, something began to shift.
After countless cycles of growth and adaptation, one of life ' s branches began using tools, and that branch grew into us.
We went on to produce a mesmerizing variety of tools, at first slowly and then with astonishing speed, we went from stone axes and fire to language, writing and eventually industrial technologies.
One invention unleashed a thousand more.
And in time, we became homo technologicus.
Around 80 years ago, another new branch of technology began.
With \textbf{the} invention of \textbf{computers}, we quickly jumped from \textbf{the} first mainframes and transistors to today ' s \textbf{smartphones} and virtual-reality headsets.
Information, knowledge, communication, computation.
In this revolution, creation has exploded like never before.
And now a new wave is upon us. \textbf{Artificial} intelligence.
These waves of history are clearly speeding up, as each one is amplified and accelerated by \textbf{the} last.
And if you look back, it ' s clear that we are in \textbf{the} fastest and most consequential wave ever. \textbf{The} journeys of humanity and technology are now deeply intertwined.
In just 18 months, over a billion people have used large language models.
We ' ve witnessed one landmark event after another.
Just a few years ago, people said that AI would never be creative.
And yet AI now feels like an endless river of creativity, making poetry and images and music and video that stretch \textbf{the} imagination.
People said it would never be empathetic.
And yet today, millions of people \textbf{enjoy} meaningful conversations with AIs, talking about their hopes and dreams and helping them work through difficult emotional challenges.
AIs can now drive cars, manage energy grids and even invent new molecules.
Just a few years ago, each of these was impossible.
And all of this is turbocharged by spiraling exponentials of data and computation.
Last year, Inflection 2. 5, our last model, used five billion times more computation than \textbf{the} DeepMind AI that beat \textbf{the} old-school Atari games just over 10 years ago.
That ' s nine orders of magnitude more computation. 10x per year, every year for almost a decade.
Over \textbf{the} same time, \textbf{the} size of these models has grown from first tens of millions of parameters to then billions of parameters, and very soon, tens of trillions of parameters.
If someone did nothing but read 24 hours a day for their entire life, they ' d consume eight billion words.
And of course, that ' s a lot of words.
But today, \textbf{the} most advanced AIs consume more than eight trillion words in a single month of training.
And all of this is set to continue. \textbf{The} long arc of technological history is now in an extraordinary new phase.
So what does this mean in practice?
Well, just as \textbf{the} internet gave us \textbf{the} browser and \textbf{the} \textbf{smartphone} gave us apps, \textbf{the} cloud-based supercomputer is ushering in a new era of ubiquitous AIs.
Everything will soon be represented by a conversational interface.
Or, to put it another way, a personal AI.
And these AIs will be infinitely knowledgeable, and soon they ' ll be factually accurate and reliable.
They ' ll have near-perfect IQ.
They ' ll also have exceptional EQ.
They ' ll be kind, supportive, empathetic.
These elements on their own would be transformational.
Just imagine if everybody had a personalized tutor in their \textbf{pocket} and access to low-cost medical advice.
A lawyer and a doctor, a business strategist and coach -- all in your \textbf{pocket} 24 hours a day.
But things really start to change when they develop what I call AQ, their"actions quotient."This is their ability to actually get stuff done in \textbf{the} digital and physical world.
And before long, it won ' t just be people that have AIs.
Strange as it may sound, every organization, from small business to nonprofit to national government, each will have their own.
Every town, building and \textbf{object} will be represented by a unique interactive persona.
And these won ' t just be mechanistic assistants.
They ' ll be companions, confidants, colleagues, friends and partners, as varied and unique as we all are.
At this point, AIs will convincingly imitate humans at most tasks.
And we ' ll feel this at \textbf{the} most intimate of scales.
An AI organizing a community get-together for an elderly neighbor.
A sympathetic expert helping you make sense of a difficult diagnosis.
But we ' ll also feel it at \textbf{the} largest scales.
Accelerating scientific discovery, autonomous cars on \textbf{the} roads, drones in \textbf{the} skies.
They ' ll both order \textbf{the} takeout and run \textbf{the} power station.
They ' ll interact with us and, of course, with each other.
They ' ll speak every language, take in every pattern of sensor data, sights, sounds, streams and streams of information, far surpassing what any one of us could consume in a thousand lifetimes.
So what is this?
What are these AIs?
If we are to prioritize safety above all else, to ensure that this new wave always serves and amplifies humanity, then we need to find \textbf{the} right metaphors for what this might become.
For years, we in \textbf{the} AI community, and I specifically, have had a tendency to refer to this as just tools.
But that doesn ' t really capture what ' s actually happening here.
AIs are clearly more dynamic, more ambiguous, more integrated and more emergent than mere tools, which are entirely subject to human control.
So to contain this wave, to put human agency at its center and to mitigate \textbf{the} inevitable unintended consequences that are likely to arise, we should start to think about them as we might a new kind of digital species.
Now it ' s just an analogy, it ' s not a literal description, and it ' s not perfect.
For a start, they clearly aren ' t biological in any traditional sense, but just pause for a moment and really think about what they already do.
They communicate in our languages.
They see what we see.
They consume unimaginably large amounts of information.
They have memory.
They have personality.
They have creativity.
They can even reason to some extent and formulate rudimentary plans.
They can act autonomously if we allow them.
And they do all this at levels of sophistication that is far beyond anything that we ' ve ever known from a mere tool.
And so saying AI is mainly about \textbf{the} math or \textbf{the} code is like saying we humans are mainly about \textbf{carbon} and water.
It ' s true, but it completely misses \textbf{the} point.
And yes, I get it, this is a super arresting thought but I honestly think this frame helps sharpen our focus on \textbf{the} critical issues.
What are \textbf{the} risks?
What are \textbf{the} boundaries that we need to impose?
What kind of AI do we want to build or allow to be built?
This is a story that ' s still unfolding.
Nothing should be accepted as a given.
We all must choose what we create.
What AIs we bring into \textbf{the} world, or not.
These are \textbf{the} questions for all of us here today, and all of us alive at this moment.
For me, \textbf{the} benefits of this technology are stunningly obvious, and they inspire my life ' s work every single day.
But quite frankly, they ' ll speak for themselves.
Over \textbf{the} years, I ' ve never shied away from highlighting risks and talking about downsides.
Thinking in this way helps us focus on \textbf{the} huge challenges that lie ahead for all of us.
But let ' s be clear.
There is no path to progress where we leave technology behind. \textbf{The} prize for all of civilization is immense.
We need solutions in health care and education, to our climate crisis.
And if AI delivers just a \textbf{fraction} of its potential, \textbf{the} next decade is going to be \textbf{the} most productive in human history.
Here ' s another way to think about it.
In \textbf{the} past, \textbf{unlocking} economic growth often came with huge downsides. \textbf{The} economy expanded as people discovered new \textbf{continents} and opened up new frontiers.
But they colonized populations at \textbf{the} same time.
We built factories, but they were grim and dangerous places to work.
We struck \textbf{oil}, but we \textbf{polluted} \textbf{the} \textbf{planet}.
Now because we are still designing and building AI, we have \textbf{the} potential and opportunity to do it better, radically better.
And today, we ' re not discovering a new \textbf{continent} and plundering its resources.
We ' re building one from scratch.
Sometimes people say that data or chips are \textbf{the} 21st century ' s new \textbf{oil}, but that ' s totally \textbf{the} wrong image.
AI is to \textbf{the} mind what nuclear \textbf{fusion} is to energy.
Limitless, abundant, world-changing.
And AI really is different, and that means we have to think about it creatively and honestly.
We have to push our analogies and our metaphors to \textbf{the} very limits to be able to grapple with what ' s coming.
Because this is not just another invention.
AI is itself an infinite inventor.
And yes, this is exciting and promising and concerning and intriguing all at once.
To be quite honest, it ' s pretty surreal.
But step back, see it on \textbf{the} long view of glacial time, and these really are \textbf{the} very most appropriate metaphors that we have today.
Since \textbf{the} beginning of life on Earth, we ' ve been \textbf{evolving}, changing and then creating everything around us in our human world today.
And AI isn ' t something outside of this story.
In fact, it ' s \textbf{the} very opposite.
It ' s \textbf{the} whole of everything that we have created, distilled down into something that we can all interact with and benefit from.
It ' s a reflection of humanity across time, and in this sense, it isn ' t a new species at all.
This is where \textbf{the} metaphors end.
Here ' s what I ' ll tell Caspian next time he asks.
AI isn ' t separate.
AI isn ' t even in some senses, new.
AI is us.
It ' s all of us.
And this is perhaps \textbf{the} most promising and vital thing of all that even a six-year-old can get a sense for.
As we build out AI, we can and must reflect all that is good, all that we love, all that is special about humanity : our empathy, our kindness, our curiosity and our creativity.
This, I would argue, is \textbf{the} greatest challenge of \textbf{the} 21st century, but also \textbf{the} most wonderful, inspiring and hopeful opportunity for all of us.
Thank you.
Chris Anderson : Thank you Mustafa.
It ' s an amazing vision and a super powerful metaphor.
You ' re in an amazing position right now.
I mean, you were connected at \textbf{the} hip to \textbf{the} amazing work happening at OpenAI.
You ' re going to have resources made available, there are reports of these giant new data centers, 100 billion dollars invested and so forth.
And a new species can emerge from it.
I mean, in your book, you did, as well as painting an incredible optimistic vision, you were super eloquent on \textbf{the} dangers of AI.
And I ' m just curious, from \textbf{the} view that you have now, what is it that most keeps you up at night?
Mustafa Suleyman : I think \textbf{the} great risk is that we get stuck in what I call \textbf{the} pessimism aversion trap.
You know, we have to have \textbf{the} courage to confront \textbf{the} potential of dark scenarios in order to get \textbf{the} most out of all \textbf{the} benefits that we see.
So \textbf{the} good news is that if you look at \textbf{the} last two or three years, there have been very, very few downsides, right?
It ' s very hard to say explicitly what harm an LLM has caused.
But that doesn ' t mean that that ' s what \textbf{the} trajectory is going to be over \textbf{the} next 10 years.
So I think if you pay attention to a few specific capabilities, take for example, autonomy.
Autonomy is very obviously a threshold over which we increase risk in our society.
And it ' s something that we should step towards very, very closely. \textbf{The} other would be something like recursive self-improvement.
If you allow \textbf{the} model to independently self-improve, update its own code, explore an environment without oversight, and, you know, without a human in control to change how it operates, that would obviously be more dangerous.
But I think that we ' re still some way away from that.
I think it ' s still a good five to 10 years before we have to really confront that.
But it ' s time to start talking about it now.
CA : A digital species, unlike any biological species, can replicate not in nine months, but in nine nanoseconds, and produce an indefinite number of \textbf{copies} of itself, all of which have more power than we have in many ways.
I mean, \textbf{the} possibility for unintended consequences seems pretty immense.
And isn ' t it true that if a problem happens, it could happen in an hour?
MS : No.
That is really not true.
I think there ' s no evidence to suggest that.
And I think that, you know, that ' s often referred to as \textbf{the}"intelligence \textbf{explosion}."And I think it is a theoretical, hypothetical maybe that we ' re all kind of curious to explore, but there ' s no evidence that we ' re anywhere near anything like that.
And I think it ' s very important that we choose our words super carefully.
Because you ' re right, that ' s one of \textbf{the} weaknesses of \textbf{the} species framing, that we will design \textbf{the} capability for self-replication into it if people choose to do that.
And I would actually argue that we should not, that would be one of \textbf{the} dangerous capabilities that we should step back from, right?
So there ' s no chance that this will"emerge"accidentally.
I really think that ' s a very low \textbf{probability}.
It will happen if engineers deliberately design those capabilities in.
And if they don ' t take enough efforts to deliberately design them out.
And so this is \textbf{the} point of being explicit and transparent about trying to introduce safety by design very early on.
CA : Thank you, your vision of humanity injecting into this new thing \textbf{the} best parts of ourselves, avoiding all those weird, biological, freaky, horrible tendencies that we can have in certain circumstances, I mean, that is a very inspiring vision.
And thank you so much for coming here and sharing it at TED.
Thank you, good luck.

% matched lemmas: artificial, carbon, computer, continent, copy, enjoy, evolve, explosion, fraction, fusion, learning, object, oil, planet, pocket, pollute, probability, smartphone, the, unlock
\end{textsample}
